\documentclass[11pt,a4paper]{beamer}
\usetheme{Madrid}
\usepackage[utf8]{inputenc}
\usepackage{amsmath}
\usepackage{amsfonts}
\usepackage{amssymb}
\usepackage{graphicx}
\usepackage{tikz}
\usepackage{booktabs}
\usepackage{array}
\usepackage{colortbl}
\usepackage{multicol}

% ============================================================
% PROFESSIONAL FONT SETTINGS (SMALL FOR DETAILED SLIDES)
% ============================================================
\usefonttheme{professionalfonts}

\setbeamerfont{normal text}{size=\tiny}
\setbeamerfont{itemize/enumerate body}{size=\tiny}
\setbeamerfont{itemize/enumerate subbody}{size=\tiny}
\setbeamerfont{block body}{size=\tiny}
\setbeamerfont{block title}{size=\scriptsize}
\setbeamerfont{frametitle}{size=\small}
\setbeamerfont{title}{size=\small}
\setbeamerfont{subtitle}{size=\footnotesize}
\setbeamerfont{author}{size=\tiny}
\setbeamerfont{date}{size=\tiny}

\setlength{\leftmargini}{0.3em}
\setlength{\leftmarginii}{0.3em}
\setlength{\labelsep}{0.1em}
\setlength{\itemsep}{0.02em}
\setlength{\parskip}{0.02em}

\title[CAMS Exam Preparation]{CAMS Exam Preparation}
\subtitle{High Net Worth Individuals - Hard Topics - Deep Dive}
\author{Private Banking | Trusts | PEPs | Offshore Risks}
\date{}

\setbeamercolor{title}{fg=white,bg=blue!70!black}
\setbeamercolor{frametitle}{fg=white,bg=blue!70!black}
\setbeamertemplate{navigation symbols}{}

\begin{document}
	
	% TITLE SLIDE
	\begin{frame}
		\titlepage
		\centering
		\tiny 30 Topics | How Questions Are Framed | Right Answers | Red Flags
	\end{frame}
	
	% ============================================================
	% SLIDE 1 - CONFLICT OF INTEREST IN PRIVATE BANKING
	% ============================================================
	\begin{frame}{1. Conflict of Interest in Private Banking}
		\tiny
		\noindent\textbf{How exam questions frame this topic:} The scenario typically describes a relationship manager who is compensated based on Assets Under Management (AUM) or revenue. The RM has a client who is a foreign PEP from a high-risk jurisdiction wanting to transfer large sums through complex offshore structures. The RM pushes to onboard quickly without Enhanced Due Diligence, telling compliance not to delay with questions. The compliance officer notices the RM has already exceeded bonus targets.
		
		\vspace{0.1cm}
		\noindent\textbf{What the right answer is:} The primary inherent financial crime risk in private banking illustrated is the conflict of interest where the RM prioritizes revenue over AML compliance. Performance compensation structures create inherent money laundering risk because RMs may overlook warning signs to maintain close relationships and generate fees.
		
		\vspace{0.1cm}
		\noindent\textbf{Common trap answers to avoid:} Answers that blame technology or process failures — such as “lack of customer identification technology” or “insufficient transaction monitoring systems” — are incorrect because even perfect technology can be overridden by an RM who fails to escalate suspicious behavior. The problem is human conflict of interest, not systems.
		
		\vspace{0.1cm}
		\noindent\textbf{Red flags to identify in the scenario:} The RM already exceeding annual bonus targets yet pushing for more. The RM telling compliance “Don't delay this with unnecessary questions.” The client being a foreign PEP from a high-risk jurisdiction. The complex structure involving shell companies and trusts. The RM focusing on \$500,000 in annual fees rather than AML concerns. Each of these elements points to revenue overriding compliance.
		
		\vspace{0.1cm}
		\noindent\textbf{How to answer correctly:} Look for the answer choice that identifies the structural conflict between compensation and compliance. The correct answer will mention “conflict of interest” and “revenue over compliance.” Eliminate any answer that blames technology or processes. Remember that regulators examine compensation structures as a key vulnerability.
	\end{frame}
	
	% ============================================================
	% SLIDE 2 - COMPLEX OWNERSHIP STRUCTURES
	% ============================================================
	\begin{frame}{2. Complex Ownership Structures (Trusts + Offshore)}
		\tiny
		\noindent\textbf{How exam questions frame this topic:} The scenario describes a high net worth client with a multi-layered structure — trust in Cook Islands, BVI holding company, Cyprus investment fund, Swiss bank account. The client refuses to provide documentation about trust beneficiaries or source of original funds, claiming privacy protected by these jurisdictions. The RM wants to onboard quickly because the client promises additional assets.
		
		\vspace{0.1cm}
		\noindent\textbf{What the right answer is:} The primary AML risks are that the complex ownership structure obscures the ultimate beneficial owner, multiple offshore jurisdictions with secrecy laws are involved, and the client refuses to disclose source of funds. Trusts can be the last layer of secrecy. Red flags include multiple offshore jurisdictions, refusal to provide documentation, and pressure to onboard quickly.
		
		\vspace{0.1cm}
		\noindent\textbf{Common trap answers to avoid:} Answers that say “no risks; trusts are legitimate wealth management tools” ignore that legitimacy depends on transparency. Legitimate trusts provide documentation. Answers that isolate only one element — “only the Swiss bank account is a risk” or “only the BVI company is a risk” — miss the cumulative risk of multiple opaque layers.
		
		\vspace{0.1cm}
		\noindent\textbf{Red flags to identify in the scenario:} Client refuses to provide trust beneficiary documentation. Client refuses to disclose source of original \$25 million. Client claims “my privacy is protected by these jurisdictions” — an admission that secrecy is intentional. Multiple offshore jurisdictions with bank secrecy laws. RM pushing to onboard despite these refusals.
		
		\vspace{0.1cm}
		\noindent\textbf{How to answer correctly:} Select the answer that lists multiple risk factors together — complex structure, multiple offshore jurisdictions, secrecy laws, and refusal to disclose. The correct answer will be comprehensive, not limited to a single element. Remember that each layer reduces transparency and increases risk.
	\end{frame}
	
	% ============================================================
	% SLIDE 3 - SOURCE OF WEALTH VS SOURCE OF FUNDS
	% ============================================================
	\begin{frame}{3. Source of Wealth vs Source of Funds}
		\tiny
		\noindent\textbf{How exam questions frame this topic:} The scenario provides client information about Source of Wealth (inheritance from family real estate business) and Source of Funds for a specific deposit (sale of cryptocurrency). The bank can verify the inheritance through a will but cannot verify the crypto sale because the client used a privacy coin and an unregulated offshore exchange. The client refuses additional documentation.
		
		\vspace{0.1cm}
		\noindent\textbf{What the right answer is:} Source of Wealth is how the client accumulated total wealth over time (the inheritance of \$25 million). Source of Funds is the specific origin of money for this particular transaction (the \$30 million crypto sale). The bank must verify BOTH independently. The unverifiable crypto sale requires Enhanced Due Diligence and potential SAR filing.
		
		\vspace{0.1cm}
		\noindent\textbf{Common trap answers to avoid:} “Source of Wealth and Source of Funds are the same; no distinction” — incorrect because they serve different purposes. “Only Source of Wealth matters for HNW clients” — incorrect because the specific deposit could be from illegal activity even if total wealth is legitimate. “Only Source of Funds matters for this deposit” — incorrect because understanding total wealth context is also required.
		
		\vspace{0.1cm}
		\noindent\textbf{Red flags to identify in the scenario:} Client used privacy coin (Monero) which makes blockchain tracing impossible. Client used unregulated offshore exchange with no KYC. Client refuses to provide additional documentation about the crypto transaction. The father's real estate business was previously investigated for money laundering (even though no charges were filed, this adds reputational risk).
		
		\vspace{0.1cm}
		\noindent\textbf{How to answer correctly:} Look for the answer that clearly distinguishes between SoW and SoF and then explains that both require independent verification. The correct answer will identify the unverifiable crypto sale as requiring EDD and potential SAR. Remember: SoW explains “how you became wealthy overall.” SoF explains “where this specific money came from.”
	\end{frame}
	
	% ============================================================
	% SLIDE 4 - TRUST STRUCTURE RED FLAGS
	% ============================================================
	\begin{frame}{4. Trust Structure Red Flags}
		\tiny
		\noindent\textbf{How exam questions frame this topic:} The scenario describes an existing trust account with a deceased settlor who was a foreign PEP. The trustee is the settlor's son (through a family office company). Trust activity includes wires to high-risk jurisdictions, purchase of a yacht titled in trust name, loan to the son's personal account with no repayment terms, and crypto purchases through an unregulated VASP. The son/trustee refuses to explain.
		
		\vspace{0.1cm}
		\noindent\textbf{What the right answer is:} The red flags include a PEP settlor with no EDD applied (PEP status applies regardless of death when concerns exist about source of funds), self-dealing transactions where the son as trustee benefits himself, unexplained transfers to high-risk jurisdiction, crypto purchases through unregulated VASP, and the trustee's refusal to explain. The bank should have applied EDD to the foreign PEP settlor.
		
		\vspace{0.1cm}
		\noindent\textbf{Common trap answers to avoid:} “No red flags; trusts can make any lawful transactions” — ignores that pattern and context matter. “Only the yacht purchase is a red flag” or “Only the cryptocurrency purchase is a red flag” — both are too narrow; multiple red flags exist. The trap argument that “the settlor is deceased; PEP status no longer matters” is incorrect.
		
		\vspace{0.1cm}
		\noindent\textbf{Red flags to identify in the scenario:} PEP settlor with no EDD applied. Self-dealing (trustee benefiting himself). Loan to son's personal account with no repayment terms. Unexplained wire to high-risk jurisdiction. Crypto through unregulated VASP. Trustee refusal to explain purpose of transactions. The principle “once a PEP, always a PEP” applies for risk assessment.
		
		\vspace{0.1cm}
		\noindent\textbf{How to answer correctly:} Select the answer that lists multiple red flags including PEP settlor with no EDD, self-dealing, unexplained transfers, crypto through unregulated VASP, and refusal to explain. Remember that self-dealing — a trustee benefiting personally from trust assets — is a serious red flag that suggests the trust is being used as a personal wallet rather than for legitimate wealth management.
	\end{frame}
	
	% ============================================================
	% SLIDE 5 - OFFSHORE FINANCIAL CENTER RED FLAGS
	% ============================================================
	\begin{frame}{5. Offshore Financial Center Red Flags}
		\tiny
		\noindent\textbf{How exam questions frame this topic:} The scenario describes a client with accounts in Cayman Islands, BVI, and Bermuda. The client's business operates entirely in one country with no offshore presence. Account activity includes round-tripping transactions, rapid transfers between offshore accounts with no economic purpose, addition of a Panama shell company as signatory, and the client is a domestic PEP.
		
		\vspace{0.1cm}
		\noindent\textbf{What the right answer is:} The red flags present include complex ownership structures, round tripping (funds moving in and out with no economic purpose), rapid asset transfers between offshore entities, use of shell companies, and presence of a PEP. Each red flag requires Enhanced Due Diligence.
		
		\vspace{0.1cm}
		\noindent\textbf{Common trap answers to avoid:} “No red flags; OFCs are legitimate for tax planning” — ignores that the pattern of activity matters, not just the existence of offshore accounts. “Only the PEP status is a red flag” or “Only the round tripping is a red flag” — both ignore the cumulative risk of multiple indicators.
		
		\vspace{0.1cm}
		\noindent\textbf{Red flags to identify in the scenario:} Round tripping — funds move Cayman to BVI to Bermuda and back to Cayman within 30 days. Rapid transfers between accounts with no economic purpose — the client's business has no operations in any offshore jurisdiction. Shell company in Panama added as additional signatory. Client is a domestic PEP. No employees or operations in offshore jurisdictions despite holding \$45 million there.
		
		\vspace{0.1cm}
		\noindent\textbf{How to answer correctly:} Look for the answer that identifies multiple red flags together. The correct answer will mention complex structures, round tripping, rapid transfers, shell companies, and PEP involvement. Remember that legitimate tax planning does not require round tripping or rapid purposeless transfers. Economic purpose is the test.
	\end{frame}
	
	% ============================================================
	% SLIDE 6 - PRIVATE INVESTMENT COMPANIES
	% ============================================================
	\begin{frame}{6. Private Investment Companies (PICs)}
		\tiny
		\noindent\textbf{How exam questions frame this topic:} The scenario describes a client who establishes a Private Investment Company in the BVI, owned by a trust where the client is both settlor and beneficiary. The client refuses to provide information about the PIC's underlying assets or investment strategy, claiming “PICs are private.” The PIC receives wires from three unrelated offshore companies and sends funds to a crypto exchange in a jurisdiction with weak AML controls.
		
		\vspace{0.1cm}
		\noindent\textbf{What the right answer is:} Private Investment Companies are considered higher risk in private banking because they reduce transparency of the ultimate beneficial owner and can be used to obscure the origin and destination of funds. The opaque nature and complex structures make it difficult to trace fund origin and destination.
		
		\vspace{0.1cm}
		\noindent\textbf{Common trap answers to avoid:} “PICs are illegal in most jurisdictions” — false; they are legal but higher risk. “PICs cannot hold bank accounts” — false; they can and do. “PICs are only used for illegal purposes” — false; legitimate PICs exist, but the structure is higher risk. The trap that “PICs are legitimate investment vehicles; no additional risk” ignores that the risk comes from opacity and refusal to disclose.
		
		\vspace{0.1cm}
		\noindent\textbf{Red flags to identify in the scenario:} Client refuses to disclose underlying assets or investment strategy. Client says “PICs are private. You don't need to know.” The PIC receives wires from three offshore companies with no apparent connection to the client. The PIC sends funds to a crypto exchange in a weak AML jurisdiction. The ownership structure (trust owning PIC, client as both settlor and beneficiary) adds opacity.
		
		\vspace{0.1cm}
		\noindent\textbf{How to answer correctly:} Select the answer that explains that PICs reduce transparency and obscure origin and destination of funds. The correct answer will focus on the opacity risk. Remember that legitimate PIC clients provide investment policy statements and quarterly holdings reports. Refusal to disclose underlying assets is grounds for EDD or rejection.
	\end{frame}
	
	% ============================================================
	% SLIDE 7 - TRUST DUE DILIGENCE REQUIREMENTS
	% ============================================================
	\begin{frame}{7. Trust Due Diligence Requirements (FATF)}
		\tiny
		\noindent\textbf{How exam questions frame this topic:} The scenario describes a trust with a deceased settlor, professional trustee (law firm), beneficiaries described only as “the children and grandchildren” with no names provided, and a protector. The law firm trustee refuses to disclose beneficiary names, citing beneficiary privacy and the discretionary nature of the trust. The bank's KYC team requests names and the law firm refuses.
		
		\vspace{0.1cm}
		\noindent\textbf{What the right answer is:} Under FATF standards, the bank must identify the settlor, all trustees, the protector (if any), AND all beneficiaries, even if discretionary. For discretionary trusts where beneficiaries are not yet determined, the bank must understand the class of beneficiaries and obtain information about known beneficiaries. Refusal to provide beneficiary information is grounds to reject the account.
		
		\vspace{0.1cm}
		\noindent\textbf{Common trap answers to avoid:} “No due diligence needed; trusts are private arrangements” — completely false under FATF. “Only the settlor and trustees need to be identified” — insufficient; beneficiaries must also be identified. “Only the law firm trustee needs to be identified” — ignores the need to identify beneficial owners.
		
		\vspace{0.1cm}
		\noindent\textbf{Red flags to identify in the scenario:} The law firm trustee cites “beneficiary privacy” — not a valid reason to withhold information from the bank. The trust is discretionary — FATF requires identification of the class of beneficiaries and known beneficiaries anyway. The trust has existed for 15 years — ample time for the trustee to have identified beneficiaries.
		
		\vspace{0.1cm}
		\noindent\textbf{How to answer correctly:} Look for the answer that requires identification of settlor, all trustees, protector, AND all beneficiaries. Remember that FATF makes no exception for discretionary trusts. If a trustee refuses to provide beneficiary names, the bank cannot establish the beneficial owner and must reject the account.
	\end{frame}
	
	% ============================================================
	% SLIDE 8 - SPECIAL PURPOSE VEHICLE RISKS
	% ============================================================
	\begin{frame}{8. Special Purpose Vehicles (SPVs)}
		\tiny
		\noindent\textbf{How exam questions frame this topic:} The scenario describes a client who establishes an SPV in Delaware to hold a \$12 million painting, uses the SPV as collateral for a loan, uses loan proceeds to buy cryptocurrency, creates a second SPV in Luxembourg to hold the crypto, refuses to explain business purpose beyond “asset protection,” and has created seven SPVs in three years each dissolved within 12-18 months.
		
		\vspace{0.1cm}
		\noindent\textbf{What the right answer is:} The financial crime risks associated with the client's use of SPVs include complex opaque structures to disguise true beneficial ownership, rapid creation and dissolution of SPVs, use of SPVs to obscure source of funds for cryptocurrency, and no clear legitimate business purpose.
		
		\vspace{0.1cm}
		\noindent\textbf{Common trap answers to avoid:} “No risks; SPVs are standard for single assets” — ignores that pattern matters. A single SPV for a single asset is normal. Seven SPVs in three years is not normal. “Only the loan transaction is a risk” or “Only the cryptocurrency purchase is a risk” — both ignore the cumulative risk of rapid creation and dissolution.
		
		\vspace{0.1cm}
		\noindent\textbf{Red flags to identify in the scenario:} Client has created seven SPVs in three years. Each SPV dissolved within 12-18 months after transferring assets. Client uses SPV as collateral for loan to buy crypto. Client creates second SPV to hold crypto. Client refuses to explain business purpose beyond vague “asset protection.” Rapid creation and dissolution suggests moving assets to avoid detection.
		
		\vspace{0.1cm}
		\noindent\textbf{How to answer correctly:} Look for the answer that mentions rapid creation and dissolution of SPVs, opaque structures, and no clear legitimate business purpose. Remember that legitimate SPVs have a clear, single purpose and are dissolved when that purpose is completed — but the pattern of many SPVs in rapid succession is suspicious.
	\end{frame}
	
	% ============================================================
	% SLIDE 9 - SOVEREIGN WEALTH FUND RISKS
	% ============================================================
	\begin{frame}{9. Sovereign Wealth Funds (SWFs)}
		\tiny
		\noindent\textbf{How exam questions frame this topic:} The scenario describes an SWF from a country with high corruption perception index ranking, board members including PEPs (finance minister, central bank governor) and two individuals with unverifiable backgrounds. The SWF refuses to provide source of funds, claiming exemption from AML as a sovereign entity. The country is on FATF Grey List.
		
		\vspace{0.1cm}
		\noindent\textbf{What the right answer is:} SWFs are not automatically exempt from AML requirements. The bank must apply Enhanced Due Diligence, verify beneficial owners (including PEP board members), assess source of funds, and consider Grey List status. The SWF's claim of exemption is incorrect.
		
		\vspace{0.1cm}
		\noindent\textbf{Common trap answers to avoid:} “No obligations; SWFs are exempt from AML” — false; FATF Recommendation 22 explicitly requires CDD for SWFs. “Only need to verify the SWF's registration” — insufficient; source of funds and beneficial owners must be verified. “Only need to screen board members against sanctions” — incomplete; EDD requires more.
		
		\vspace{0.1cm}
		\noindent\textbf{Red flags to identify in the scenario:} SWF from high-corruption country (low CPI ranking). Board includes PEPs (finance minister, central bank governor). Two board members with unverifiable backgrounds. SWF refuses to provide source of funds. SWF claims exemption incorrectly. Country on FATF Grey List for 18 months.
		
		\vspace{0.1cm}
		\noindent\textbf{How to answer correctly:} Look for the answer that states SWFs are not automatically exempt and requires EDD, PEP verification, source of funds assessment, and Grey List consideration. Remember that sovereign immunity does not waive AML obligations. Treat SWFs like any other institutional client.
	\end{frame}
	
	% ============================================================
	% SLIDE 10 - HIGH VALUE ASSET RED FLAGS
	% ============================================================
	\begin{frame}{10. High Value Asset Red Flags}
		\tiny
		\noindent\textbf{How exam questions frame this topic:} The scenario describes a client purchasing multiple high-value assets over six months: a painting paid by wire from Seychelles, a diamond necklace paid in three cash installments, a yacht registered to a Marshall Islands shell company, and gold bars through a dealer with no AML program. Client's declared annual income is \$3 million but purchases total \$85 million. Client sells each asset within 12-18 months at a loss.
		
		\vspace{0.1cm}
		\noindent\textbf{What the right answer is:} The red flags indicating money laundering include funds from high-risk jurisdiction (Seychelles), cash structuring (\$4 million installments), shell company ownership, dealer with no AML, income inconsistency (\$3 million income vs \$85 million purchases), and rapid resale at loss. Criminals accept losses to clean money.
		
		\vspace{0.1cm}
		\noindent\textbf{Common trap answers to avoid:} “No red flags; HNW individuals buy luxury assets regularly” — ignores the pattern of inconsistent income, cash structuring, and rapid resale at loss. “Only the cash payments are a red flag” or “Only the shell company registration is a red flag” — both ignore the cumulative risk of multiple indicators.
		
		\vspace{0.1cm}
		\noindent\textbf{Red flags to identify in the scenario:} Funds from secrecy jurisdiction (Seychelles). Cash structuring — breaking \$12 million into three \$4 million installments. Shell company ownership of yacht. Unregulated dealer with no AML program. Income inconsistency — \$3 million income cannot explain \$85 million purchases. Rapid resale at loss — selling within 12-18 months at a loss is not rational investment behavior.
		
		\vspace{0.1cm}
		\noindent\textbf{How to answer correctly:} Select the answer that lists multiple red flags including funds from high-risk jurisdiction, cash structuring, shell company, unregulated dealer, income inconsistency, and rapid resale at loss. Remember that rapid resale at a consistent loss is the strongest indicator — criminals accept losses to clean money.
	\end{frame}
	
	% ============================================================
	% SLIDE 11 - FAMILY OFFICE RISKS
	% ============================================================
	\begin{frame}{11. Family Office Risks}
		\tiny
		\noindent\textbf{How exam questions frame this topic:} The scenario describes a Single Family Office managing \$500 million for a family whose patriarch was a former government official (PEP) who left office 8 years ago. The family office is registered in Delaware but managed from Switzerland. It refuses to identify an external investment advisor. Transactions include wires to 15 countries including two FATF Grey List jurisdictions, and unexplained incoming wires from a BVI company.
		
		\vspace{0.1cm}
		\noindent\textbf{What the right answer is:} The primary AML risks include PEP involvement (patriarch — once a PEP always a PEP for risk purposes), opaque multi-jurisdictional structure (Delaware registration, Swiss management), unidentified third parties (external advisor), transactions to Grey List jurisdictions, and unexplained incoming wires from BVI shell company jurisdiction.
		
		\vspace{0.1cm}
		\noindent\textbf{Common trap answers to avoid:} “No risks; family offices are low risk” — ignores that family offices can be used for ML. “Only the PEP status is a risk” or “Only the Grey List transactions are a risk” — both ignore the cumulative risk. The trap that “patriarch left office 8 years ago, no longer PEP” is incorrect.
		
		\vspace{0.1cm}
		\noindent\textbf{Red flags to identify in the scenario:} PEP patriarch — once a PEP always a PEP. Multi-jurisdictional structure reduces transparency. Unidentified external advisor — family office refuses to provide information. Transactions to FATF Grey List jurisdictions. Unexplained incoming wires from BVI shell company with no connection to known family business.
		
		\vspace{0.1cm}
		\noindent\textbf{How to answer correctly:} Look for the answer that mentions PEP involvement (with note that once a PEP always a PEP), opaque multi-jurisdictional structure, unidentified third parties, transactions to Grey List jurisdictions, and unexplained incoming wires. Remember that family offices should be transparent about all service providers.
	\end{frame}
	
	% ============================================================
	% SLIDE 12 - INSURANCE PRODUCTS MONEY LAUNDERING
	% ============================================================
	\begin{frame}{12. Insurance Products Money Laundering Technique}
		\tiny
		\noindent\textbf{How exam questions frame this topic:} The scenario describes a client purchasing a single-premium variable life insurance policy for \$10 million using funds from an offshore account. Three months later, the client takes a policy loan of \$8 million (80\% of cash value) and wires proceeds to a real estate developer. Six months later, the client surrenders the policy and receives \$1.5 million as a check. This same pattern has been repeated with three different insurance companies in two years.
		
		\vspace{0.1cm}
		\noindent\textbf{What the right answer is:} The money laundering technique being attempted is policy loan layering. The pattern is: (1) purchase policy with illicit funds, (2) take policy loan to receive clean funds (layering), (3) surrender policy for additional clean funds (integration). Repeating the pattern with multiple insurers indicates systematic laundering.
		
		\vspace{0.1cm}
		\noindent\textbf{Common trap answers to avoid:} “No laundering; insurance is a legitimate investment” — ignores the pattern of purchase, loan, and early surrender across multiple insurers. “Only the surrender proceeds are laundered” or “Only the loan proceeds are laundered” — both miss that both the loan and surrender produce clean funds from the original dirty premium.
		
		\vspace{0.1cm}
		\noindent\textbf{Red flags to identify in the scenario:} Single-premium policy purchased with funds from offshore account. Policy loan taken very soon after purchase (three months). Loan proceeds wired to third party (real estate developer). Early surrender of policy (six months — much earlier than typical). Pattern repeated with three different insurance companies in two years.
		
		\vspace{0.1cm}
		\noindent\textbf{How to answer correctly:} Select the answer that identifies “policy loan layering” and explains the three-step pattern. Remember that the repeating pattern across multiple insurers is key — a single policy might be legitimate, but systematic repetition is laundering.
	\end{frame}
	
	% ============================================================
	% SLIDE 13 - CRYPTOCURRENCY AND HNW CLIENTS
	% ============================================================
	\begin{frame}{13. Cryptocurrency \& HNW Clients}
		\tiny
		\noindent\textbf{How exam questions frame this topic:} The scenario describes a client with declared net worth of \$50 million from real estate development who suddenly engages in significant cryptocurrency activity: purchasing \$8 million in Bitcoin through a VASP with weak KYC, transferring to DeFi protocol earning 12\% APY, converting profits through decentralized exchange, sending to regulated Swiss VASP, and cashing out \$9.5 million. The client's real estate business has declining revenues for three years. The client refuses to provide wallet addresses or transaction history.
		
		\vspace{0.1cm}
		\noindent\textbf{What the right answer is:} The red flags present include disconnect between declining business performance and large crypto purchases, use of VASP with weak KYC, DeFi protocol with no transparency, refusal to provide wallet history, and unexplained rapid profit. EDD is required.
		
		\vspace{0.1cm}
		\noindent\textbf{Common trap answers to avoid:} “No red flags; HNW clients invest in crypto” — ignores the disconnect between declining business income and large crypto purchases. “Only the VASP choice is a red flag” or “Only the profit amount is a red flag” — both ignore the cumulative risk.
		
		\vspace{0.1cm}
		\noindent\textbf{Red flags to identify in the scenario:} Declining business revenues for three years but \$8 million crypto purchase — where did the money come from? VASP with weak KYC. DeFi protocol with no transparency — bank cannot verify transactions. Client refuses to provide wallet addresses or transaction history. Unexplained rapid profit from DeFi.
		
		\vspace{0.1cm}
		\noindent\textbf{How to answer correctly:} Look for the answer that mentions disconnect between business performance and crypto activity, weak KYC VASP, DeFi opacity, refusal to provide wallet history, and unexplained profit. Remember that declining income plus large purchases suggests undisclosed source of funds.
	\end{frame}
	
	% ============================================================
	% SLIDE 14 - ART AND COLLECTIBLES LAUNDERING
	% ============================================================
	\begin{frame}{14. Art and Collectibles Laundering}
		\tiny
		\noindent\textbf{How exam questions frame this topic:} The scenario describes a client purchasing \$50 million in fine art over three years. The pattern: art bought through London auction house, shipped to Geneva freeport, client never takes possession, art sold 6-12 months later at 20-30\% loss, sale proceeds deposited to different bank account than purchase account. Client creates single-purpose LLCs for each purchase/sale then dissolves them. Client refuses to explain strategy.
		
		\vspace{0.1cm}
		\noindent\textbf{What the right answer is:} The indicators that art is being used for money laundering rather than investment include freeport storage (no physical possession), rapid resale at loss, use of single-purpose LLCs, refusal to explain, and different bank accounts for purchase and sale.
		
		\vspace{0.1cm}
		\noindent\textbf{Common trap answers to avoid:} “No indicators; art investment is legitimate” — ignores the pattern of rapid resale at loss and freeport storage. “Only the freeport storage is suspicious” or “Only the resale at loss is suspicious” — both ignore the cumulative pattern.
		
		\vspace{0.1cm}
		\noindent\textbf{Red flags to identify in the scenario:} Freeport storage — client never takes physical possession. Rapid resale at 20-30\% loss within 6-12 months. Single-purpose LLCs created for each transaction then dissolved. Different bank accounts for purchase and sale. Client refuses to explain strategy, claiming “art is a private matter.”
		
		\vspace{0.1cm}
		\noindent\textbf{How to answer correctly:} Select the answer that lists freeport storage, rapid resale at loss, single-purpose LLCs, refusal to explain, and different accounts. Remember that legitimate art collectors buy art they enjoy, display it, and hold for years. Criminals flip art quickly at a loss.
	\end{frame}
	
	% ============================================================
	% SLIDE 15 - SECURED LOAN RISKS
	% ============================================================
	\begin{frame}{15. Secured Loan Risks}
		\tiny
		\noindent\textbf{How exam questions frame this topic:} The scenario describes a client with a \$15 million portfolio of publicly traded stocks and bonds who requests a secured loan of \$12 million for “business expansion.” Three days after receiving the loan, the client wires \$10 million to a company in a high-risk jurisdiction with no connection to the client's business. The client then sells the collateral portfolio and repays the loan early. The business has no record of receiving the \$10 million.
		
		\vspace{0.1cm}
		\noindent\textbf{What the right answer is:} The secured loan structure can be used for money laundering by using illicit funds as collateral (the portfolio may be purchased with dirty money). The loan proceeds are clean and can be moved freely. Early repayment after selling collateral creates a legitimate-appearing source of funds.
		
		\vspace{0.1cm}
		\noindent\textbf{Common trap answers to avoid:} “Loans cannot be used for money laundering” — false; this is a known technique. “Only the early repayment is suspicious” or “Only the wire to high-risk jurisdiction is suspicious” — both ignore that the entire structure enables laundering.
		
		\vspace{0.1cm}
		\noindent\textbf{Red flags to identify in the scenario:} Loan purpose stated as “business expansion” but business has no record of receiving funds. Wire to high-risk jurisdiction with no connection to client's business. Early repayment of loan after selling collateral. The original portfolio used as collateral may have been purchased with illicit funds.
		
		\vspace{0.1cm}
		\noindent\textbf{How to answer correctly:} Look for the answer that explains that illicit funds used as collateral generate clean loan proceeds, and early repayment creates a legitimate paper trail. Remember that secured loans are not low risk for ML — the collateral itself may be dirty.
	\end{frame}
	
	% ============================================================
	% SLIDE 16 - MULTI-JURISDICTION TAX EVASION
	% ============================================================
	\begin{frame}{16. Multi-Jurisdiction Tax Evasion}
		\tiny
		\noindent\textbf{How exam questions frame this topic:} The scenario describes a client who is citizen of Country A, tax resident of Country B, operates business in Country C, and has bank accounts in Countries D, E, and F. The client moves funds weekly between accounts with no business purpose. Tax returns show \$2 million annual income but total assets are \$200 million. Client refuses to sign CRS self-certification form and requests bank not report interest to any tax authority.
		
		\vspace{0.1cm}
		\noindent\textbf{What the right answer is:} Tax evasion is a predicate offense for money laundering in most jurisdictions. The bank must file CRS reports despite client refusal. Unexplained cross-border movements and income inconsistency require EDD and potential SAR.
		
		\vspace{0.1cm}
		\noindent\textbf{Common trap answers to avoid:} “No obligations; tax evasion is not a predicate offense” — false in most jurisdictions. “Only need to file CRS if client signs form” — false; CRS reporting is mandatory regardless. “Bank cannot report without client consent” — false; CRS requires reporting without consent.
		
		\vspace{0.1cm}
		\noindent\textbf{Red flags to identify in the scenario:} Client moves funds weekly between accounts in three countries with no business purpose. Income inconsistency — \$2 million income but \$200 million assets. Refusal to sign CRS self-certification form. Client requests bank not report interest to any tax authority.
		
		\vspace{0.1cm}
		\noindent\textbf{How to answer correctly:} Select the answer that states tax evasion is a predicate offense, CRS reporting is mandatory despite client refusal, and unexplained movements plus income inconsistency require EDD and SAR. Remember that refusal to provide CRS form is itself a red flag.
	\end{frame}
	
	% ============================================================
	% SLIDE 17 - SHELL COMPANY NETWORKS
	% ============================================================
	\begin{frame}{17. Shell Company Networks}
		\tiny
		\noindent\textbf{How exam questions frame this topic:} The scenario describes a client controlling 25 shell companies across eight jurisdictions (BVI, Cayman, Delaware, Panama, Seychelles, Bahamas, Belize, Mauritius). The client's declared business is international consulting with \$5 million annual revenue. Transaction monitoring shows circular fund movement among all 25 companies with total annual flow of \$500 million. Client refuses to provide financial statements and recently added a Cook Islands trust as owner of the entire network.
		
		\vspace{0.1cm}
		\noindent\textbf{What the right answer is:} The primary money laundering risk is the opaque ownership structure obscuring fund movement. Circular transactions indicate layering. The revenue-to-flow ratio (\$5 million revenue vs \$500 million flow) is wildly inconsistent. The trust overlay adds a final secrecy layer.
		
		\vspace{0.1cm}
		\noindent\textbf{Common trap answers to avoid:} “No risk; shell companies are legal business structures” — ignores the massive discrepancy between revenue and flow. “Only the circular transactions are a risk” or “Only the revenue inconsistency is a risk” — both miss that the combination creates the strongest red flag.
		
		\vspace{0.1cm}
		\noindent\textbf{Red flags to identify in the scenario:} 25 shell companies across 8 secrecy jurisdictions. Declared revenue \$5 million but transaction flow \$500 million — 100 times higher. Circular fund movement among all companies. Client refuses to provide financial statements. Trust overlay added as owner with client as settlor and beneficiary.
		
		\vspace{0.1cm}
		\noindent\textbf{How to answer correctly:} Look for the answer that highlights the revenue-to-flow discrepancy as the strongest red flag, along with circular transactions and trust overlay. Remember that 100x flow-to-revenue ratio is mathematically impossible for a legitimate business.
	\end{frame}
	
	% ============================================================
	% SLIDE 18 - ROUND TRIPPING
	% ============================================================
	\begin{frame}{18. Round Tripping Detection}
		\tiny
		\noindent\textbf{How exam questions frame this topic:} The scenario describes a client with accounts at two banks. Over 12 months: Bank A sends \$1 million to Bank B; Bank B sends \$980,000 back to Bank A five days later. This repeats with different amounts: \$2M sent, \$1.95M returned; \$3M sent, \$2.92M returned; \$5M sent, \$4.88M returned. Each round trip loses 2-5\%. Total moved \$50 million with \$2 million lost. Client cannot explain purpose.
		
		\vspace{0.1cm}
		\noindent\textbf{What the right answer is:} The most likely purpose is money laundering layering — creating complex transaction history to obscure fund origin. Criminals are willing to accept losses (the 2-5\% is the laundering cost). There is no legitimate business purpose, and the client's inability to explain requires a SAR.
		
		\vspace{0.1cm}
		\noindent\textbf{Common trap answers to avoid:} “Currency hedging” — hedging does not involve consistent losses. “Tax avoidance” — this pattern does not create tax benefits. “Legitimate cash management” — legitimate cash management does not involve sending money out and receiving less back repeatedly.
		
		\vspace{0.1cm}
		\noindent\textbf{Red flags to identify in the scenario:} Consistent loss percentage (2-5\%) on each round trip. Same pattern repeating with different amounts. Client cannot explain purpose. Business is manufacturing with no cross-border operations — no legitimate reason for this pattern. Total \$2 million lost over 12 months — accepting losses indicates laundering.
		
		\vspace{0.1cm}
		\noindent\textbf{How to answer correctly:} Select the answer that identifies money laundering layering and notes that criminals accept losses to clean money. Remember that a consistent loss percentage is a fee — that's the laundering cost.
	\end{frame}
	
	% ============================================================
	% SLIDE 19 - FOREIGN PEP ENHANCED DUE DILIGENCE
	% ============================================================
	\begin{frame}{19. Foreign PEP Enhanced Due Diligence}
		\tiny
		\noindent\textbf{How exam questions frame this topic:} The scenario describes a foreign PEP who was a former prime minister who left office 5 years ago. The PEP wants to deposit \$25 million. Source of wealth: pension (\$500,000 annually) plus consulting fees (\$2 million annually) — total \$2.5 million per year. Over 5 years, maximum legitimate income is \$12.5 million before taxes. The PEP cannot explain the gap. The PEP refuses to provide tax returns from before government service. The PEP's country ranks 145/180 on Corruption Perceptions Index.
		
		\vspace{0.1cm}
		\noindent\textbf{What the right answer is:} FATF requires EDD for foreign PEPs regardless of time since leaving office. The bank must verify source of wealth and source of funds, obtain senior management approval, and conduct enhanced monitoring. The mathematical inconsistency (\$25 million deposit vs \$12.5 million maximum possible legitimate income) requires a SAR if not resolved.
		
		\vspace{0.1cm}
		\noindent\textbf{Common trap answers to avoid:} “No EDD needed; PEP left office 5 years ago” — false; FATF requires EDD regardless of time. “Only need senior management approval” — insufficient; source verification also required. “Only need source of funds verification” — insufficient; source of wealth also required.
		
		\vspace{0.1cm}
		\noindent\textbf{Red flags to identify in the scenario:} PEP left office 5 years ago but FATF still requires EDD. Mathematical inconsistency — \$25 million cannot be explained by \$2.5 million annual income over 5 years. PEP claims “savings before becoming prime minister” but refuses to provide documentation. Country has high corruption perception index (145/180).
		
		\vspace{0.1cm}
		\noindent\textbf{How to answer correctly:} Look for the answer that states EDD is required regardless of time since leaving office, includes source verification and senior management approval, and notes that mathematical inconsistency requires SAR. Remember: do the math. If income cannot explain wealth, file SAR.
	\end{frame}
	
	% ============================================================
	% SLIDE 20 - DOMESTIC PEP RISK ASSESSMENT
	% ============================================================
	\begin{frame}{20. Domestic PEP Risk Assessment}
		\tiny
		\noindent\textbf{How exam questions frame this topic:} The scenario describes a domestic PEP — senior government official in the same country where the bank operates. Declared income is \$300,000 from government salary. The official wants to deposit \$2 million from “family inheritance.” Bank verifies father's estate was \$1.5 million. The additional \$500,000 cannot be explained. The official's son won three government contracts worth \$800,000 where the official was involved in awarding them.
		
		\vspace{0.1cm}
		\noindent\textbf{What the right answer is:} FATF recommends a risk-based assessment for domestic PEPs, not automatic EDD. However, red flags in this case — unexplained wealth (\$500,000 gap) and son's government contracts with father's involvement — require EDD regardless of PEP type.
		
		\vspace{0.1cm}
		\noindent\textbf{Common trap answers to avoid:} “Domestic PEPs are always low risk” — false; risk depends on circumstances. “Domestic PEPs require same EDD as foreign PEPs” — not under FATF; risk-based approach applies. “No EDD needed for domestic PEPs” — false when red flags exist.
		
		\vspace{0.1cm}
		\noindent\textbf{Red flags to identify in the scenario:} Unexplained wealth — \$500,000 gap after accounting for verified inheritance. Son's company won government contracts where father was involved in awarding — potential conflict of interest and corruption. Country has strong anti-corruption laws but weak enforcement.
		
		\vspace{0.1cm}
		\noindent\textbf{How to answer correctly:} Look for the answer that explains FATF's risk-based approach for domestic PEPs but notes that red flags (unexplained wealth and son's contracts) require EDD. Remember that red flags override any presumption of low risk.
	\end{frame}
	
	% ============================================================
	% SLIDE 21 - INTERNATIONAL ORGANIZATION PEPs
	% ============================================================
	\begin{frame}{21. International Organization PEPs (e.g., World Bank)}
		\tiny
		\noindent\textbf{How exam questions frame this topic:} The scenario describes a senior executive at the World Bank (international organization). Declared annual income \$400,000. Client wants to deposit \$3 million from “savings over 20 years.” Bank statements show monthly deposits of \$10,000-15,000 for 20 years, totaling approximately \$3 million. However, deposits stop during the client's 5 years of World Bank employment and resume after leaving. The client cannot explain why deposits stopped. Client's country of origin is a high-risk jurisdiction.
		
		\vspace{0.1cm}
		\noindent\textbf{What the right answer is:} FATF recommends risk-based assessment for international organization PEPs. The deposit pattern inconsistency — no deposits during World Bank employment when income was highest — is a red flag requiring EDD. The client cannot explain why.
		
		\vspace{0.1cm}
		\noindent\textbf{Common trap answers to avoid:} “No EDD needed; international organization PEPs are low risk” — false; risk-based assessment applies. “Same as foreign PEP: automatic EDD required” — not under FATF; risk-based applies. “No obligations; World Bank has its own AML” — false; the bank still has obligations.
		
		\vspace{0.1cm}
		\noindent\textbf{Red flags to identify in the scenario:} Deposit pattern inconsistency — deposits stopped during high-income World Bank years and resumed after leaving. Client cannot explain why. Client's country of origin is a high-risk jurisdiction. The pattern suggests deposits before and after World Bank may have been fabricated.
		
		\vspace{0.1cm}
		\noindent\textbf{How to answer correctly:} Select the answer that notes risk-based assessment applies and that inconsistent deposit patterns require EDD. Remember: wealth should accumulate faster during high-income years, not stop.
	\end{frame}
	
	% ============================================================
	% SLIDE 22 - IMMEDIATE FAMILY OF PEPs
	% ============================================================
	\begin{frame}{22. Immediate Family of PEPs}
		\tiny
		\noindent\textbf{How exam questions frame this topic:} The scenario describes the spouse of a foreign PEP (current cabinet minister). The spouse has a consulting firm with \$500,000 annual revenue. The spouse wants to deposit \$10 million described as “gift from husband.” The husband's declared government salary is \$150,000 annually with no other declared wealth. The family's lifestyle includes \$5 million home, private jets, and luxury vacations. The spouse refuses to provide documentation about the gift.
		
		\vspace{0.1cm}
		\noindent\textbf{What the right answer is:} FATF definition includes immediate family members (spouse, children, parents, siblings). The spouse is a PEP by association. EDD is required. The \$10 million gift from a \$150,000 salary is mathematically impossible — SAR required.
		
		\vspace{0.1cm}
		\noindent\textbf{Common trap answers to avoid:} “Family members are not PEPs; no EDD needed” — false; FATF explicitly includes immediate family. “Only the PEP spouse requires EDD, not the family member” — false; the family member is also a PEP. “No EDD needed for gifts” — false; source of gift must be verified.
		
		\vspace{0.1cm}
		\noindent\textbf{Red flags to identify in the scenario:} Spouse of foreign PEP (current cabinet minister). \$10 million gift from husband with \$150,000 salary — mathematically impossible. Family lifestyle (\$5 million home, private jets, luxury vacations) inconsistent with declared income. Spouse refuses to provide gift documentation. Spouse claims “my husband's finances are private.”
		
		\vspace{0.1cm}
		\noindent\textbf{How to answer correctly:} Look for the answer that states immediate family are PEPs under FATF, EDD is required, and the mathematical impossibility requires SAR. Remember: if the math doesn't work, file SAR.
	\end{frame}
	
	% ============================================================
	% SLIDE 23 - CLOSE ASSOCIATES OF PEPs
	% ============================================================
	\begin{frame}{23. Close Associates of PEPs}
		\tiny
		\noindent\textbf{How exam questions frame this topic:} The scenario describes a long-time friend and business partner of a foreign PEP (former intelligence agency director). The client and PEP co-own a company that won 15 government contracts worth \$200 million over 8 years. The PEP's ownership is hidden through a trust. The client is the public-facing owner. The client refuses to identify other owners. Open source research shows allegations that the PEP used his position to steer contracts to the company.
		
		\vspace{0.1cm}
		\noindent\textbf{What the right answer is:} Close associates are included in the PEP definition under FATF. Hidden ownership structure, allegations of contract steering, and refusal to identify owners all require EDD and potential SAR.
		
		\vspace{0.1cm}
		\noindent\textbf{Common trap answers to avoid:} “No special treatment; only the PEP matters” — false; close associates are included. “Only need to screen the client against PEP lists” — insufficient; EDD required. “No action unless charges are filed” — false; allegations and red flags require action.
		
		\vspace{0.1cm}
		\noindent\textbf{Red flags to identify in the scenario:} Client is close associate (business partner) of foreign PEP. PEP's ownership hidden through trust. Client refuses to identify other owners. Company won government contracts where PEP had influence. News articles allege contract steering.
		
		\vspace{0.1cm}
		\noindent\textbf{How to answer correctly:} Select the answer that states close associates are included in PEP definition, hidden ownership and allegations require EDD, and potential SAR should be considered. Remember that close associates are often used as frontmen for corrupt PEPs.
	\end{frame}
	
	% ============================================================
	% SLIDE 24 - CORRESPONDENT BANKING RISKS
	% ============================================================
	\begin{frame}{24. Correspondent Banking Risks (FATF Rec. 13)}
		\tiny
		\noindent\textbf{How exam questions frame this topic:} The scenario describes a private bank in Country A with a correspondent relationship with a private bank in Country B (offshore financial center). Country B bank has 500 private banking clients including 50 foreign PEPs. Country B refuses to provide client names citing bank secrecy laws. Country B processes \$10 billion annually in wires. Country B has been cited for AML deficiencies twice in 3 years and has no independent AML audit in 4 years.
		
		\vspace{0.1cm}
		\noindent\textbf{What the right answer is:} Under FATF Recommendation 13, the bank must terminate or restrict the relationship. The respondent bank has AML deficiencies, refuses to identify underlying clients (including PEPs), operates in an offshore center, and has no independent audit.
		
		\vspace{0.1cm}
		\noindent\textbf{Common trap answers to avoid:} “No obligation; correspondent banking is standard” — false; EDD is required. “Only need to file annual report” — insufficient. “Only need to monitor transactions” — insufficient when respondent refuses UBO information.
		
		\vspace{0.1cm}
		\noindent\textbf{Red flags to identify in the scenario:} Respondent refuses to identify underlying clients including PEPs — cites bank secrecy laws. Respondent has AML deficiencies cited twice in 3 years. No independent AML audit in 4 years. Respondent operates in offshore financial center.
		
		\vspace{0.1cm}
		\noindent\textbf{How to answer correctly:} Look for the answer that states termination or restriction is required. Remember that without UBO identification, the correspondent bank is blind to PEP exposure and sanctions risk.
	\end{frame}
	
	% ============================================================
	% SLIDE 25 - NESTED ACCOUNTS
	% ============================================================
	\begin{frame}{25. Nested Accounts in Correspondent Banking}
		\tiny
		\noindent\textbf{How exam questions frame this topic:} The scenario describes a private bank with a correspondent account for Respondent Bank R. Bank R allows its private banking clients to use the correspondent account directly (nested accounts). The private bank discovers one of Bank R's clients is a foreign PEP from a sanctioned country. The private bank never approved this nesting. Bank R refuses to terminate the PEP client's access, stating “the client is important to our business.” Investigation reveals 30\% of transactions are for Bank R's private banking clients, not Bank R itself.
		
		\vspace{0.1cm}
		\noindent\textbf{What the right answer is:} Unauthorized nesting is a serious risk. The private bank must file a SAR, demand termination of nesting, and consider terminating the correspondent relationship. The correspondent bank has no KYC information for underlying clients including PEPs from sanctioned countries.
		
		\vspace{0.1cm}
		\noindent\textbf{Common trap answers to avoid:} “No action; nesting is standard industry practice” — false when unauthorized. “Only need to file SAR for the PEP client” — insufficient; the nesting arrangement itself is the problem. “Only need to monitor transactions more closely” — insufficient when respondent refuses to terminate nesting.
		
		\vspace{0.1cm}
		\noindent\textbf{Red flags to identify in the scenario:} Unauthorized nesting — private bank never approved. PEP from sanctioned country using nested account. Bank R refuses to terminate PEP client's access. 30\% of transactions are for underlying private banking clients, not Bank R itself. Private bank has no KYC on any underlying clients.
		
		\vspace{0.1cm}
		\noindent\textbf{How to answer correctly:} Select the answer that requires SAR filing, demand termination of nesting, and consideration of terminating the correspondent relationship. Remember: unauthorized nesting is never acceptable.
	\end{frame}
	
	% ============================================================
	% SLIDE 26 - EXITING HIGH-RISK CLIENTS
	% ============================================================
	\begin{frame}{26. Exiting High-Risk Clients}
		\tiny
		\noindent\textbf{How exam questions frame this topic:} The scenario describes a bank that decided to exit a high-risk client — a foreign PEP with unexplained wealth of \$30 million who refuses to provide source of funds documentation. The client has outstanding loans of \$8 million secured by real estate and a \$10 million irrevocable letter of credit. The client threatens to sue for breach of contract. Legal team advises litigation risk. RM urges keeping client for \$500,000 annual fees.
		
		\vspace{0.1cm}
		\noindent\textbf{What the right answer is:} The bank's obligation is to exit despite legal and commercial pressures. Exit strategy must consider legal obligations but cannot retain client solely for revenue. Manage loan and LC separately (allow to run off or call if contract permits). File SAR. Document exit decision.
		
		\vspace{0.1cm}
		\noindent\textbf{Common trap answers to avoid:} “Keep the client; cannot exit due to loans and LCs” — false; exit is possible with careful management. “Only exit after loans are repaid” — delays exit unnecessarily; loans can be allowed to run off. “Never exit PEP clients; manage with EDD only” — false; exit is required when risk cannot be mitigated.
		
		\vspace{0.1cm}
		\noindent\textbf{Red flags to identify in the scenario:} Client is foreign PEP with unexplained wealth. Client refuses to provide source of funds documentation. Outstanding loans and LC are used as arguments to keep client. RM cites \$500,000 annual fees. Client threatens litigation.
		
		\vspace{0.1cm}
		\noindent\textbf{How to answer correctly:} Look for the answer that states exit is required despite loans and LC, manage legal obligations separately, file SAR, and document decision. Remember: revenue is never a valid reason to retain a high-risk client.
	\end{frame}
	
	% ============================================================
	% SLIDE 27 - CULTURE OF COMPLIANCE
	% ============================================================
	\begin{frame}{27. Culture of Compliance}
		\tiny
		\noindent\textbf{How exam questions frame this topic:} The scenario describes internal audit findings: the bank's 5 largest clients (all foreign PEPs) contribute 40\% of private banking revenue. RMs for these 5 clients have never filed a SAR despite significant red flags. Compliance recommended terminating 3 of the 5 clients. Senior management rejected all 3 termination recommendations citing “commercial sensitivity.” RM bonus structure is 100\% revenue-based with no risk adjustment. Board has not reviewed any AML metrics in 2 years.
		
		\vspace{0.1cm}
		\noindent\textbf{What the right answer is:} This indicates a weak culture of compliance. Revenue overrides AML. Senior management overrides compliance. No Board oversight. Compensation structure encourages risk-taking. This is a textbook failure of compliance culture.
		
		\vspace{0.1cm}
		\noindent\textbf{Common trap answers to avoid:} “Strong culture; compliance recommendations are considered” — false; they were rejected. “Only the bonus structure is problematic” — insufficient; multiple failures exist. “Only the lack of Board review is problematic” — insufficient; senior management overrides are more serious.
		
		\vspace{0.1cm}
		\noindent\textbf{Red flags to identify in the scenario:} Largest clients are foreign PEPs contributing 40\% of revenue — concentration risk. RMs never filed SARs despite red flags. Compliance recommended termination; senior management rejected. Bonus structure 100\% revenue-based. Board no AML review in 2 years.
		
		\vspace{0.1cm}
		\noindent\textbf{How to answer correctly:} Select the answer that identifies weak culture of compliance, revenue overriding AML, senior management overrides, no Board oversight, and problematic compensation. Remember: senior management overriding compliance is the most serious failure.
	\end{frame}
	
	% ============================================================
	% SLIDE 28 - ESTATE PLANNING TRUST ABUSES
	% ============================================================
	\begin{frame}{28. Estate Planning Trust Abuses}
		\tiny
		\noindent\textbf{How exam questions frame this topic:} The scenario describes a trust for estate planning: settlor age 75, trustees (settlor and son), beneficiaries (settlor life interest, then son and grandchildren). Client funds trust with \$40 million. Three months later, settlor resigns as trustee, leaving son as sole trustee. Six months later, trust sells \$30 million in assets and distributes proceeds to son personally (not as trustee). Son purchases \$25 million yacht and \$5 million crypto. Son moves to country with no extradition. Father claims no knowledge. Son has no independent source of wealth.
		
		\vspace{0.1cm}
		\noindent\textbf{What the right answer is:} Red flags include rapid resignation of settlor as trustee, distributions to son personally (not as trustee), son has no source of wealth, son's relocation to non-extradition country, and father's claim of no knowledge. This suggests potential asset dissipation or fraud requiring SAR.
		
		\vspace{0.1cm}
		\noindent\textbf{Common trap answers to avoid:} “No red flags; estate planning is legitimate” — ignores the pattern. “Only the yacht purchase is a red flag” or “Only the cryptocurrency purchase is a red flag” — both ignore the more serious issue of distributions to son personally.
		
		\vspace{0.1cm}
		\noindent\textbf{Red flags to identify in the scenario:} Settlor resigns as trustee only three months after funding trust. Distributions to son personally, not in trustee capacity. Son has no independent source of wealth. Son moves to non-extradition country. Father claims no knowledge. Assets used for personal luxury (yacht) and crypto.
		
		\vspace{0.1cm}
		\noindent\textbf{How to answer correctly:} Look for the answer that mentions rapid resignation, distributions personally, son's lack of wealth, relocation to non-extradition country, and father's claim of no knowledge. Remember: distributions to an individual not in trustee capacity is the key red flag.
	\end{frame}
	
	% ============================================================
	% SLIDE 29 - SUCCESSION PLANNING RED FLAGS
	% ============================================================
	\begin{frame}{29. Succession Planning Red Flags}
		\tiny
		\noindent\textbf{How exam questions frame this topic:} The scenario describes an 80-year-old client with \$250 million net worth creating trusts for three children (\$80 million each) and a foundation (\$10 million) in Cook Islands. After client's death, all three trusts transfer assets to new trusts in Marshall Islands with Child A as trustee of new trusts. Child B moves to country with no financial disclosure laws. Child C purchases \$30 million in crypto through unregulated VASPs. Foundation's \$10 million wired to Seychelles with no charitable activity.
		
		\vspace{0.1cm}
		\noindent\textbf{What the right answer is:} Red flags include migration to secrecy jurisdictions (Cook Islands to Marshall Islands), self-dealing (Child A as trustee), relocation to non-disclosure country, crypto purchases through unregulated VASPs, and foundation funds diverted from charitable purpose. This suggests possible asset hiding or money laundering.
		
		\vspace{0.1cm}
		\noindent\textbf{Common trap answers to avoid:} “No red flags; HNW clients can structure succession as they wish” — ignores the pattern of moving to increasingly secretive jurisdictions. “Only the foundation diversion is a red flag” or “Only the crypto purchases are a red flag” — both ignore multiple other indicators.
		
		\vspace{0.1cm}
		\noindent\textbf{Red flags to identify in the scenario:} Migration from Cook Islands to Marshall Islands (increasing secrecy). Child A becomes trustee of new trusts (self-dealing). Child B moves to non-disclosure country. Child C buys \$30 million crypto through unregulated VASPs. Foundation funds to Seychelles with no charitable activity.
		
		\vspace{0.1cm}
		\noindent\textbf{How to answer correctly:} Select the answer that lists migration to secrecy jurisdictions, self-dealing, relocation to non-disclosure country, crypto through unregulated VASPs, and foundation diversion. Remember: foundation funds diverted from charitable purpose is loss of charitable purpose and potential ML.
	\end{frame}
	
	% ============================================================
	% SLIDE 30 - OFFSHORE TRUST HIGHEST RISK FEATURES
	% ============================================================
	\begin{frame}{30. Offshore Trust Highest Risk Features}
		\tiny
		\noindent\textbf{How exam questions frame this topic:} The scenario describes an offshore trust in the Bahamas with features: settlor from no-AML jurisdiction, bank-affiliated trustee, Swiss lawyer as protector, beneficiaries as “family members” with no names, flee clause allowing asset movement if threatened, letter of wishes updated annually, assets include cash, art in Geneva freeport, and crypto with unregulated VASP. Trust refuses to provide letter of wishes or identify beneficiaries. Settlor has created and dissolved 6 similar trusts in 10 years, each lasting 18-24 months.
		
		\vspace{0.1cm}
		\noindent\textbf{What the right answer is:} The highest risk features are anonymous beneficiaries (no CDD possible), flee clause (asset mobility), letter of wishes (secret instructions), rapid creation/dissolution of trusts, art in freeport (opaque asset), crypto with unregulated VASP, and settlor from no-AML jurisdiction. This is a high-risk structure requiring EDD and likely SAR.
		
		\vspace{0.1cm}
		\noindent\textbf{Common trap answers to avoid:} “No features present risk; offshore trusts are legitimate” — ignores that these specific features are designed to evade detection. “Only the crypto assets are a risk” or “Only the freeport art is a risk” — both ignore the cumulative risk of multiple secrecy features.
		
		\vspace{0.1cm}
		\noindent\textbf{Red flags to identify in the scenario:} Anonymous beneficiaries — no CDD possible. Flee clause — allows rapid asset movement to avoid freezing. Letter of wishes — secret instructions not reviewed by bank. Pattern of trust creation/dissolution (each 18-24 months). Art in freeport — opaque valuation and ownership. Crypto with unregulated VASP. Settlor from jurisdiction with no AML controls.
		
		\vspace{0.1cm}
		\noindent\textbf{How to answer correctly:} Look for the answer that lists multiple risk features including anonymous beneficiaries, flee clause, letter of wishes, rapid dissolution pattern, freeport art, and unregulated crypto. Remember: trusts with features designed to evade detection or freezing should be treated as high risk.
	\end{frame}
	
	% ============================================================
	% SUMMARY SLIDE
	% ============================================================
	\begin{frame}{Summary: How to Answer HNW Exam Questions}
		\tiny
		\noindent\textbf{Key principles for answering correctly:}
		
		\vspace{0.1cm}
		\noindent• When you see a relationship manager pushing to onboard quickly despite red flags, the answer is always conflict of interest from AUM-based compensation. Eliminate technology or process answers.
		
		\vspace{0.1cm}
		\noindent• When you see multiple offshore layers and refusal to disclose, the answer is that complex structures obscure UBO and require EDD. The correct answer lists multiple red flags, not just one.
		
		\vspace{0.1cm}
		\noindent• When you see Source of Wealth verified but Source of Funds unverified, the answer is that both must be independently verified. Unverifiable SoF requires EDD and SAR.
		
		\vspace{0.1cm}
		\noindent• When you see a trust with PEP settlor (even deceased) and self-dealing, the answer is that PEP status continues and self-dealing is a red flag.
		
		\vspace{0.1cm}
		\noindent• When you see round tripping or rapid transfers between offshore accounts with no economic purpose, the answer is layering requiring SAR.
		
		\vspace{0.1cm}
		\noindent• When you see rapid resale of luxury assets at a loss, the answer is money laundering — the loss is the laundering cost.
		
		\vspace{0.1cm}
		\noindent• When you see mathematical inconsistency between income and wealth, the answer is EDD and SAR — do the math.
		
		\vspace{0.1cm}
		\noindent• When you see foreign PEPs, the answer is always EDD regardless of time since leaving office.
		
		\vspace{0.1cm}
		\noindent• When you see domestic PEPs with red flags (unexplained wealth, family contracts), the answer is EDD despite risk-based approach.
		
		\vspace{0.1cm}
		\noindent• When you see unauthorized nested accounts in correspondent banking, the answer is terminate relationship.
		
		\vspace{0.1cm}
		\noindent• When you see senior management overriding compliance for revenue, the answer is culture of compliance failure.
		
		\vspace{0.1cm}
		\noindent• When in doubt, file SAR. SAR is almost always the correct action when red flags remain unexplained.
		
		\vspace{0.3cm}
		\centering
		{Remember: Identify the red flags, do the math, and choose the answer that requires action (EDD, SAR, termination).}
	\end{frame}
	
\end{document}